% =====================================================================
\chapter{Összefoglalás és további munka}
\label{ch:osszefoglalas}
% =====================================================================

\section{Összefoglalás}

Dolgozatunkban a vizsgakérdés-generálás négy gyakorlati problémájából indultunk ki: a kulcsok
helyessége nincs ellenőrizve (P1), a vizsga szerkezete és lefedettsége nincs kontrollálva (P2), a
tananyag egy harmadik félhez kerül (P3), és a módszer nagy, távoli modellekre támaszkodik (P4). Ezekre a
\Mimir{} rendszerrel válaszoltunk, amely az egyetlen promptot explicit, ellenőrizhető módszerrel váltja
fel, és egyetlen fogyasztói GPU-n, adatmegőrzés nélkül fut.

A munka eredményei a következőkben foglalhatók össze.

\begin{enumerate}[label=(\arabic*)]
\item \textbf{Formális modell.} A „jó vizsga” fogalmát mérhető, forrásalapú predikátumokra bontottuk
(formai helyesség, alátámasztottság, disztraktor-validitás, vak megoldhatóság, lefedettség,
szivárgás), a generálást korlátos optimalizálási feladatként írtuk fel, és bizonyítottuk a vizsgaterv
két szerkezeti tulajdonságát: a Bloom-kiosztás kvótatulajdonságát és a mohó slotkiosztás
lefedettségi alsó korlátját.
\item \textbf{Generálási módszer.} Kidolgoztuk és implementáltuk a tervezésvezérelt, forráshivatkozásos
csővezetéket: ablakozott szemantikus darabolás, lefedettségvezérelt vizsgaterv, slotonkénti sűrű vagy
hibrid visszakeresés, egyenkénti forráshoz kötött generálás, legolcsóbbtól a legdrágábbig rendezett
ellenőrző lánc visszacsatolással, tartalékfogalmakkal és javító hívással, valamint munkamenet-hatókörű
fogalomgráf.
\item \textbf{Adatvédő architektúra.} A mikroszolgáltatás-architektúra zéró adatmegőrzéssel,
ujjlenyomatokra épülő naplózással és kapcsolható, kizárólag helyi üzemmóddal működik; a GPU-memóriaigény
analitikus modelljét mérésekkel validáltuk: a 7B generátor és futtatókörnyezete 5,4--5,9\,GiB-ot foglal,
így 8\,GB-os kártyán is fut.
\item \textbf{Kiértékelési módszertan.} Tizenöt egytényezős kísérleti kart, 39 dokumentumos kétnyelvű
adatkészletet 326 referenciakérdéssel, modellcsaládon kívüli bírót, közös bíró--oktató rubrikát és
dokumentumszintű statisztikát tartalmazó, nyílt és reprodukálható keretrendszert hoztunk létre.
\item \textbf{Pilot eredmények.} Egyetlen dokumentumon a tervezés javította a helyi modell
alátámasztottságát, vak megoldhatóságát, disztraktorait és lefedettségét; az azonos modellel végzett
önellenőrzés a bíróval csak véletlenszintű egyezést mutatott; a visszacsatolási hurok a kulcs kérdésbe
szivárgását idézte elő, amely az automatikus mutatókat és a bírót egyaránt félrevezette; a futásról
futásra jelentkező zaj ($\approx0{,}07$) pedig megmutatta, hogy egyetlen dokumentumon a legtöbb különbség
nem értelmezhető. Ezekből a megfigyelésekből közvetlenül következtek a fő vizsgálat javításai.
\item \textbf{A fő vizsgálat eredményei.} Tizennégy karral, kilenc magyar és angol dokumentumon
(karonként hat), 1\,670 kérdésen: a helyi 7B modellel futó tervezett és ellenőrzött csővezeték érte el a
legmagasabb minőségi indexet a közös dokumentumokon ($0{,}96$), szemben a naiv RAG $0{,}64$-es, a helyi teljes
dokumentumos prompt $0{,}83$-as és a teljes dokumentumot kapó szervermodell $0{,}95$-ös értékével; a tervezés
$+0{,}24$-dal, az ellenőrzés -- a pilottal ellentétben -- mind a hat dokumentumon, $+0{,}08$-dal javított; a
más modellcsaládba tartozó ellenőrző és a fogalomgráf rontott; a szivárgási szabályok a szivárgást
0--3\,\%-ra szorították le; a naiv RAG az angol dokumentumokon a rövid szerkezeti darabok miatt összeomlott;
ugyanaz a csővezeték a szervermodellel $0{,}99$-et ért el. A Holm-korrekció után egyik összevetés sem
szignifikáns.
\end{enumerate}

A kutatási kérdésekre a két vizsgálat alapján a következő válaszokat adhatjuk. \textbf{RQ1:} a strukturált
generálás a vizsgált dokumentumokon érvényesebb vizsgát adott, mint ha ugyanannak a modellnek a dokumentumot
vagy annak visszakeresett részleteit egyetlen promptban adjuk át ($0{,}96$ vs.\ $0{,}83$ és $0{,}64$), és a
lefedettsége jobb volt a naiv RAG-énál; a teljes dokumentumos promptnál azonban nem volt jobb a lefedettsége
(\ref{h:lefed}.~hipotézis nem teljesült). \textbf{RQ2:} a tervezés 17 hívásért, az ellenőrzés további
34 hívásért javított, a vak megoldási próba a lánc hasznos eleme; a más modellcsaládba tartozó ellenőrző,
a fogalomgráf és a hibrid keresés a többletköltségéért nem javított; a naiv RAG-ot a darabolás apró
részletei is jelentősen befolyásolták. \textbf{RQ3:} a helyi 7B modell a tervezett és ellenőrzött
csővezetékben a közös dokumentumokon elérte a szervermodell naiv RAG-os ($0{,}95$ vs.\ $0{,}93$) és teljes
dokumentumos ($0{,}96$ vs.\ $0{,}95$) változatát, a pilotban mért 5,4--5,9\,GiB GPU-memóriával; ugyanazzal a
csővezetékkel a szervermodell $0{,}99$-et ért el, vizsgánként 47,7 perc alatt.

\section{További munka}

A közvetlen folytatás a fő vizsgálat kiterjesztése \az{\ref{ch:kiertekeles}}.~fejezetben rögzített
protokoll szerint: a fő karok (\arm{B-doc-L}, \arm{B-doc-S}, \arm{E0}, \arm{E1}, \arm{E2}, \arm{E4},
\arm{E4b}) futtatása mind a 39 dokumentumon, köztük a 24 magyar Wikipédia-cikken, azonos dokumentumokon
a helyi és a szerverkarokkal; a bíró validálása három oktató vak értékelésével; és a bíró kontextusablakának
bővítése a teljes dokumentumos karokra. Ezen túl a következő irányokat tartjuk ígéretesnek.

\begin{itemize}
\item \textbf{Minimális darabméret.} A naiv RAG angol dokumentumokon mért összeomlását a szerkezeti
sorokból képzett, tartalom nélküli darabok okozták; a darabolóba illesztett minimális darabméret (a rövid
darab összevonása a szomszédjával) egyszerű javítás, amelynek hatása egy újabb \arm{E0}-futással mérhető.
\item \textbf{Megbízhatóbb ellenőrző.} Az ellenőrzés javítja a vizsgát, de a kérdésenkénti jelzése a bíróval
csak gyengén egyezik, és a más modellcsaládba tartozó ellenőrző (\arm{E2x}) sem jobb. Természetes nyelvi
következtetésre (NLI) tanított kisebb osztályozó vagy a FActScore-hoz hasonló atomi
tényellenőrzés~\cite{min2023factscore} lehet a következő lépés.
\item \textbf{A megvalósult Bloom-szint ellenőrzése.} A pilot szerint a kis modell a kért kognitív
szintet nem valósítja meg megbízhatóan; egy szintosztályozó ellenőrzés az ellenőrző láncba illeszthető.
\item \textbf{Adaptív ellenőrzés.} Az ellenőrzés költségét csökkentené, ha csak a bizonytalan
kérdésekre futna (pl.\ a kulcs lexikális alátámasztottsága vagy a generátor bizonyossága alapján).
\item \textbf{Nyelvi összevetés.} Az angol dokumentumok egy részének magyar fordítása lehetővé tenné a
nyelv hatásának a forrástól független mérését.
\item \textbf{További kérdéstípusok és tanórai kipróbálás.} Az igaz--hamis és a kifejtős kérdések
kiértékelése, valamint a generált vizsgák pszichometriai elemzése (nehézségi és diszkriminációs index,
nem funkcionáló disztraktorok) valós hallgatói válaszokon.
\end{itemize}

A \Mimir{} nyílt forráskódú: a kód, a konfigurációk, az adatkészlet-manifesztek, a pilot nyers
kimenetei és a fő vizsgálat jelentése a reprodukáláshoz rendelkezésre állnak.
